\chapter{Module Architecture}
\label{ch:architecture}

\section{Module Dependency Graph}

OmniLaTeX loads its 27 library modules in a strict order to satisfy
dependencies.  The load order forms a directed acyclic graph with no
circular dependencies:

\begin{enumerate}
  \item \textbf{Core} -- \texttt{omnilatex-base.sty} is always loaded
        first by \texttt{omnilatex.cls}.  It provides fundamental macros
        and sets up the package infrastructure.

  \item \textbf{Utils} -- Utility modules loaded early because they are
        depended on by nearly everything else:
        \texttt{omnilatex-utils}, \texttt{omnilatex-colors},
        \texttt{omnilatex-themes}.

  \item \textbf{Typography} -- Font, list, typesetting, and math modules:
        \texttt{omnilatex-fonts}, \texttt{omnilatex-lists},
        \texttt{omnilatex-typesetting}, \texttt{omnilatex-math}.

  \item \textbf{Layout} -- Page geometry, KOMA integration, floats,
        boxes, document structure, presentation, and accessibility:
        \texttt{omnilatex-page}, \texttt{omnilatex-koma},
        \texttt{omnilatex-floats}, \texttt{omnilatex-boxes},
        \texttt{omnilatex-document}, \texttt{omnilatex-presentation},
        \texttt{omnilatex-accessibility}.

  \item \textbf{Graphics} -- Image support and TikZ integration:
        \texttt{omnilatex-graphics}, \texttt{omnilatex-tikz-core},
        \texttt{omnilatex-tikz-engineering}.

  \item \textbf{References} -- Hyperlinks, bibliography, citations, and
        glossary: \texttt{omnilatex-hyperref}, \texttt{omnilatex-biblio},
        \texttt{omnilatex-citations}, \texttt{omnilatex-glossary}.

  \item \textbf{Code} -- Code listing support:
        \texttt{omnilatex-listings}.

  \item \textbf{Tables} -- Table formatting:
        \texttt{omnilatex-tables}.

  \item \textbf{Language} -- Internationalisation, CJK, and RTL:
        \texttt{omnilatex-i18n}, \texttt{omnilatex-cjk},
        \texttt{omnilatex-rtl}.
\end{enumerate}

Optional modules are gated by \texttt{enable*} class options (see
Chapter~\ref{ch:class-options}).  When disabled, the corresponding
\texttt{\textbackslash RequirePackage} call is skipped entirely, reducing
load time and avoiding unnecessary package conflicts.

\section{Module Categories}

All 27 modules are organised into nine categories by functional domain:

\begin{table}[htbp]
  \centering
  \begin{tabular}{@{} l l p{5cm} @{}}
    \toprule
    \textbf{Category} & \textbf{Modules} & \textbf{Count} \\
    \midrule
    Core &
      omnilatex-base &
      1 \\[3pt]
    Utils &
      omnilatex-utils, omnilatex-colors, omnilatex-themes &
      3 \\[3pt]
    Typography &
      omnilatex-fonts, omnilatex-lists, omnilatex-typesetting,
      omnilatex-math &
      4 \\[3pt]
    Layout &
      omnilatex-page, omnilatex-koma, omnilatex-floats,
      omnilatex-boxes, omnilatex-document, omnilatex-presentation,
      omnilatex-accessibility &
      7 \\[3pt]
    Graphics &
      omnilatex-graphics, omnilatex-tikz-core,
      omnilatex-tikz-engineering &
      3 \\[3pt]
    References &
      omnilatex-hyperref, omnilatex-biblio, omnilatex-citations,
      omnilatex-glossary &
      4 \\[3pt]
    Code &
      omnilatex-listings &
      1 \\[3pt]
    Tables &
      omnilatex-tables &
      1 \\[3pt]
    Language &
      omnilatex-i18n, omnilatex-cjk, omnilatex-rtl &
      3 \\
    \midrule
    \textbf{Total} & & \textbf{27} \\
    \bottomrule
  \end{tabular}
  \caption{Module categories and their constituent packages}
\end{table}

\section{Design Principles}

\subsection{Single Responsibility}

Each module has a clearly defined scope.  \texttt{omnilatex-fonts}
handles only font selection; \texttt{omnilatex-floats} handles only
float placement; \texttt{omnilatex-i18n} handles only language and
translation setup.  Cross-cutting concerns (e.g.\ a float that uses a
custom font) are handled by the composing document, not by merging module
boundaries.

\subsection{Optional Loading}

Modules that pull in heavy dependencies are gated behind boolean class
options.  The six optional feature flags are:

\begin{table}[htbp]
  \centering
  \begin{tabular}{@{} l l l @{}}
    \toprule
    \textbf{Option} & \textbf{Default} & \textbf{Modules Enabled} \\
    \midrule
    \texttt{enablefonts} & false & omnilatex-fonts (extended) \\
    \texttt{enablegraphics} & false & omnilatex-graphics \\
    \texttt{enablemath} & true & omnilatex-math \\
    \texttt{enabletikz} & true & omnilatex-graphics, omnilatex-tikz-core \\
    \texttt{enableengineering} & true & omnilatex-tikz-engineering \\
    \texttt{enablecode} & true & omnilatex-listings \\
    \texttt{enabletables} & true & omnilatex-tables \\
    \bottomrule
  \end{tabular}
  \caption{Feature flags controlling optional module loading}
\end{table}

\subsection{KOMA Integration}

Layout modules integrate tightly with KOMA-Script rather than fighting
against it.  \texttt{omnilatex-koma.sty} provides the bridge layer that
translates OmniLaTeX's doctype system into KOMA class options
(\texttt{\textbackslash LoadClass}), while \texttt{omnilatex-page.sty}
uses KOMA's \ctanpackage{typearea} for margin calculations rather than
manual \texttt{\textbackslash geometry} calls.

\subsection{No Circular Dependencies}

The load order is strictly hierarchical: Core $\to$ Utils $\to$
Typography $\to$ Layout $\to$ Graphics $\to$ References $\to$ Code $\to$
Tables $\to$ Language.  A module in one tier may depend on modules in
earlier tiers but never on modules in the same or later tiers.  This is
enforced by the sequential \texttt{\textbackslash RequirePackage} calls
in \texttt{omnilatex.cls}.

\subsection{File Location}

All library modules reside under \texttt{lib/} with one subdirectory per
category:

\begin{minted}{text}
lib/
  core/          omnilatex-base.sty
  utils/         omnilatex-utils.sty, omnilatex-colors.sty, omnilatex-themes.sty
  typography/    omnilatex-fonts.sty, omnilatex-lists.sty, ...
  layout/        omnilatex-page.sty, omnilatex-koma.sty, ...
  graphics/      omnilatex-graphics.sty, omnilatex-tikz-core.sty, ...
  references/    omnilatex-hyperref.sty, omnilatex-biblio.sty, ...
  code/          omnilatex-listings.sty
  tables/        omnilatex-tables.sty
  language/      omnilatex-i18n.sty, omnilatex-cjk.sty, omnilatex-rtl.sty
\end{minted}
